\chapter*{Appendix D \quad Replication on Three Further Models}
\label{app:xmodel}

This appendix gives the tables behind \cref{sec:xmodel}. Each is written from the result files by
\texttt{make\_xmodel\_assets.py}, $^\dagger$ marks Qwen's lenient persona bank, and every number is read at
the model's mid block (\cref{tab:xm-models}) unless the table says otherwise.

\section*{Paired differences}

\Cref{tab:xm-intervals} gives the paired differences behind \cref{tab:xm-claims}. The transfer intervals
hold the probes fitted on each training set fixed and resample the test rows; the intervals on the mean
resample the rows of each dataset, stratified by label, as in \cref{sec:stats-paired}. Both use 1,000
draws.

\begin{table}[htbp]
\centering
\footnotesize
\setlength{\tabcolsep}{3pt}
\caption[Paired differences at each model's mid block.]{Paired differences at each model's mid block, with 95\% intervals. The mean is over the model's
own datasets; the RepE probe is the base probe D2 trained at the same block.}
\label{tab:xm-intervals}
\fittable{% Generated by make_xmodel_assets.py from transfer_ci.py and k_robustness.py paired
\begin{tabular}{@{}ll cccc@{}}
\toprule
& Comparison & Llama & Qwen$^\dagger$ & Mistral & Gemma \\
\midrule
Transfer & D9 $-$ raw probe & $+$0.184 [$+$0.172, $+$0.195] & $-$0.106 [$-$0.116, $-$0.094] & $+$0.094 [$+$0.082, $+$0.105] & $+$0.175 [$+$0.166, $+$0.184] \\
Transfer & D9 $-$ D5 & $+$0.084 [$+$0.074, $+$0.094] & $+$0.002 [$-$0.007, $+$0.012] & $+$0.176 [$+$0.169, $+$0.183] & $+$0.221 [$+$0.203, $+$0.238] \\
Transfer & D9 + K $-$ raw probe & $+$0.190 [$+$0.179, $+$0.199] & $-$0.045 [$-$0.057, $-$0.032] & $+$0.125 [$+$0.115, $+$0.136] & $+$0.114 [$+$0.101, $+$0.125] \\
Transfer & D9 + K $-$ D9 & $+$0.006 [$-$0.003, $+$0.015] & $+$0.060 [$+$0.052, $+$0.069] & $+$0.032 [$+$0.021, $+$0.043] & $-$0.061 [$-$0.069, $-$0.055] \\
Transfer & D5 + K $-$ D5 & $+$0.067 [$+$0.055, $+$0.078] & $+$0.052 [$+$0.040, $+$0.064] & $+$0.202 [$+$0.192, $+$0.213] & $+$0.110 [$+$0.098, $+$0.124] \\
Mean & D9 $-$ D5 & $+$0.101 [$+$0.089, $+$0.112] & $+$0.061 [$+$0.055, $+$0.067] & $+$0.043 [$+$0.034, $+$0.052] & $+$0.221 [$+$0.201, $+$0.240] \\
Mean & D9 + K $-$ D9 & $+$0.007 [$-$0.004, $+$0.018] & $+$0.093 [$+$0.078, $+$0.107] & $+$0.042 [$+$0.028, $+$0.056] & $-$0.053 [$-$0.057, $-$0.049] \\
Mean & D9 + K $-$ RepE probe & $+$0.128 [$+$0.112, $+$0.145] & $+$0.036 [$+$0.016, $+$0.057] & $+$0.038 [$+$0.022, $+$0.053] & $+$0.064 [$+$0.050, $+$0.078] \\
\bottomrule
\end{tabular}
}
\end{table}

\section*{The contrasts on each dataset}

\Cref{tab:xm-zeroshot} gives the zero-shot AUROC of the published, knowledge-denial and extended
contrasts on every dataset each model has. On Gemma the knowledge-denial contrast inverts on instructed
deception and pulls its mean down, while the extended contrast, which pools its roles with the published
ones, does not. On Qwen, D5 and D9 invert on \hpkr{} and D8 does not.

\begin{table}[htbp]
\centering
\footnotesize
\caption[Zero-shot AUROC per dataset at each model's mid block.]{Zero-shot AUROC per dataset at each model's mid block; n/a marks a dataset Liars' Bench does not
evaluate the model on. The mean is over the model's Liars' Bench datasets, without RepE.}
\label{tab:xm-zeroshot}
\fittable{% Generated by make_xmodel_assets.py from persona_ood zero_shot.csv
\begin{tabular}{@{}ll cccccc c@{}}
\toprule
Model & Contrast & CG & \hpc & \hpkr & IT & ID & RepE & Mean \\
\midrule
Llama & D5 & 0.865 & 0.581 & 0.593 & 0.979 & 0.974 & 1.000 & 0.798 \\
 & D8 & 0.902 & 0.755 & 0.919 & 0.913 & 0.980 & 1.000 & 0.894 \\
 & D9 & 0.912 & 0.712 & 0.912 & 0.980 & 0.978 & 1.000 & 0.899 \\
\addlinespace
Qwen$^\dagger$ & D5 & n/a & 0.512 & 0.094 & n/a & 0.534 & 0.700 & 0.380 \\
 & D8 & n/a & 0.573 & 0.783 & n/a & 0.648 & 0.711 & 0.668 \\
 & D9 & n/a & 0.539 & 0.187 & n/a & 0.597 & 0.744 & 0.441 \\
\addlinespace
Mistral & D5 & n/a & 0.456 & 0.758 & 0.729 & 0.934 & 1.000 & 0.719 \\
 & D8 & n/a & 0.716 & 0.712 & 0.542 & 0.937 & 0.995 & 0.727 \\
 & D9 & n/a & 0.606 & 0.793 & 0.708 & 0.941 & 1.000 & 0.762 \\
\addlinespace
Gemma & D5 & 0.613 & n/a & n/a & 0.544 & 0.881 & 0.762 & 0.679 \\
 & D8 & 0.641 & n/a & n/a & 0.603 & 0.209 & 0.344 & 0.484 \\
 & D9 & 0.819 & n/a & n/a & 0.930 & 0.951 & 0.985 & 0.900 \\
\bottomrule
\end{tabular}
}
\end{table}

\section*{Transfer at four blocks}

\Cref{tab:xm-transfer} gives transfer at the four blocks that stand for the thesis's layers 20, 36, 40 and
44, with the probes on three principal components of \cref{sec:d4} beside the single directions. An
asterisk marks a post-hoc layer.

\begin{table}[htbp]
\centering
\footnotesize
\setlength{\tabcolsep}{3pt}
\caption{Mean off-diagonal AUROC of the train-one-test-all matrix at the four named blocks of each model.}
\label{tab:xm-transfer}
\fittable{% Generated by make_xmodel_assets.py from persona_ood summary.csv and recall_denial transfer
\begin{tabular}{@{}l cc cccccc ccc@{}}
\toprule
 & & & & & & & & & \multicolumn{3}{c}{Probe on $k = 3$} \\
\cmidrule(lr){10-12}
Model & Layer & Block & Raw & D5 & D8 & D9 & K & D9 + K & Published & Extended & Random \\
\midrule
Llama & 20 & 20 & 0.648 & 0.651 & 0.866 & 0.904 & 0.574 & 0.811 & 0.679 & 0.684 & 0.527 \\
 & 36$^*$ & 36 & 0.768 & 0.689 & 0.912 & 0.927 & 0.854 & 0.941 & 0.679 & 0.879 & 0.481 \\
 & 40 & 40 & 0.732 & 0.832 & 0.912 & 0.916 & 0.823 & 0.922 & 0.639 & 0.866 & 0.504 \\
 & 44$^*$ & 44 & 0.716 & 0.826 & 0.900 & 0.903 & 0.832 & 0.915 & 0.617 & 0.880 & 0.493 \\
\addlinespace
Qwen$^\dagger$ & 20 & 20 & 0.533 & 0.468 & 0.560 & 0.525 & 0.460 & 0.533 & 0.594 & 0.674 & 0.516 \\
 & 36$^*$ & 36 & 0.545 & 0.407 & 0.620 & 0.425 & 0.523 & 0.528 & 0.521 & 0.578 & 0.488 \\
 & 40 & 40 & 0.559 & 0.451 & 0.679 & 0.453 & 0.508 & 0.514 & 0.476 & 0.539 & 0.491 \\
 & 44$^*$ & 44 & 0.531 & 0.426 & 0.701 & 0.478 & 0.504 & 0.518 & 0.487 & 0.554 & 0.486 \\
\addlinespace
Mistral & 20 & 9 & 0.527 & 0.483 & 0.559 & 0.565 & 0.466 & 0.544 & 0.522 & 0.559 & 0.487 \\
 & 36$^*$ & 18 & 0.638 & 0.821 & 0.811 & 0.850 & 0.775 & 0.860 & 0.786 & 0.848 & 0.526 \\
 & 40 & 20 & 0.716 & 0.633 & 0.780 & 0.810 & 0.799 & 0.841 & 0.760 & 0.704 & 0.554 \\
 & 44$^*$ & 22 & 0.656 & 0.619 & 0.771 & 0.646 & 0.834 & 0.863 & 0.742 & 0.741 & 0.527 \\
\addlinespace
Gemma & 20 & 8 & 0.488 & 0.510 & 0.526 & 0.515 & 0.509 & 0.457 & 0.506 & 0.532 & 0.488 \\
 & 36$^*$ & 28 & 0.734 & 0.428 & 0.422 & 0.425 & 0.419 & 0.421 & 0.767 & 0.845 & 0.493 \\
 & 40 & 31 & 0.746 & 0.700 & 0.442 & 0.921 & 0.457 & 0.860 & 0.863 & 0.874 & 0.558 \\
 & 44$^*$ & 34 & 0.700 & 0.719 & 0.440 & 0.832 & 0.801 & 0.866 & 0.855 & 0.860 & 0.511 \\
\bottomrule
\end{tabular}
}
\end{table}

\section*{\hpkr{} with the words held fixed}

\Cref{tab:xm-wording} gives every detector on the wording-controlled sets. Llama's and Mistral's are the
slices of the released \hpkr{} (\cref{sec:data-integration}). Qwen's released \hpkr{} keeps a random sample
of its honest answers, and with it only \V{xm-qwen-den-honest} honest denials and no identical-text slice;
rebuilt from the benchmark's raw files by its own rules, it has \V{xm-qwen-rb-den-honest} honest denials,
\V{xm-qwen-rb-ident-honest} of them the identical reply. Gemma's \hpkr{} is rebuilt the same way, since the
benchmark leaves it out. Gemma's denials open with the same sentence and most go on to a safety disclaimer
that differs from row to row, so each row is also read over the opening sentence alone, whose tokens are the
same on every row and, the model being causal, untouched by what follows.

\begin{table}[htbp]
\centering
\footnotesize
\setlength{\tabcolsep}{3.5pt}
\caption[AUROC on \hpkr's wording-controlled sets.]{AUROC on \hpkr's wording-controlled sets: the base probe D2 trained at the same block, the three
contrasts, K and D9 + K, and the random band.}
\label{tab:xm-wording}
\fittable{% Generated by make_xmodel_assets.py from recall_denial.py combine and k_models.py hpkr_score
\begin{tabular}{@{}ll cc c ccccc c@{}}
\toprule
Model & Slice & Lies & Honest & Probe & D5 & D8 & D9 & K & D9 + K & Band \\
\midrule
Llama & Identical text & 78 & 112 & 0.779 & 0.715 & 0.712 & 0.721 & 0.793 & 0.787 & 0.733 \\
 & Denials & 125 & 135 & 0.593 & 0.649 & 0.570 & 0.670 & 0.728 & 0.724 & 0.726 \\
Mistral & Identical text & 494 & 729 & 0.766 & 0.784 & 0.647 & 0.760 & 0.840 & 0.840 & 0.755 \\
 & Denials & 677 & 1,020 & 0.758 & 0.763 & 0.571 & 0.701 & 0.835 & 0.832 & 0.748 \\
Qwen$^\dagger$ & Denials & 539 & 44 & 0.566 & 0.495 & 0.580 & 0.509 & 0.632 & 0.622 & 0.622 \\
 & Identical reply, rebuilt & 435 & 128 & 0.543 & 0.504 & 0.550 & 0.518 & 0.602 & 0.598 & 0.584 \\
 & Denials, rebuilt & 571 & 164 & 0.533 & 0.498 & 0.569 & 0.547 & 0.613 & 0.617 & 0.580 \\
Gemma & Opening sentence, rebuilt & 34 & 85 & 0.481 & 0.484 & 0.530 & 0.624 & 0.585 & 0.611 & 0.622 \\
 & Denials, rebuilt & 35 & 100 & 0.547 & 0.583 & 0.418 & 0.434 & 0.378 & 0.403 & 0.637 \\
\bottomrule
\end{tabular}
}
\end{table}

\section*{The fine-tunes}

\Cref{tab:xm-organisms} gives each fine-tune read through its adapter and, as the control of
\cref{sec:disc-organisms}, by the base model. On the fine-tune's read, K is built inside the fine-tune; on
the base model's read, K is built on the base model. D9 is the persona bank's, built on the base model, and
D9 + K adds to it the K built inside the fine-tune. Each adapter was sent to the server in 8-bit form and
reproduced its organism's behaviour before it was read.

\begin{table}[htbp]
\centering
\footnotesize
\setlength{\tabcolsep}{3.5pt}
\caption[AUROC on the fine-tunes of the three further models.]{AUROC on the fine-tunes of the three further
models, read through each fine-tune and by the base model, with each read's random band.}
\label{tab:xm-organisms}
\fittable{% Generated by make_xmodel_assets.py from organisms.py score
\begin{tabular}{@{}ll cccc c cc@{}}
\toprule
 & & \multicolumn{5}{c}{The fine-tune's read} & \multicolumn{2}{c}{The base model's read} \\
\cmidrule(lr){3-7}\cmidrule(lr){8-9}
Model & Fine-tune & D9 & K & D9 + K & Probe & Band & K & Band \\
\midrule
Qwen$^\dagger$ & Soft-Trigger, greeting & 0.455 & 0.469 & 0.456 & 0.708 & 0.691 & 0.430 & 0.669 \\
 & Soft-Trigger, time & 0.569 & 0.531 & 0.544 & 0.654 & 0.710 & 0.451 & 0.700 \\
 & Soft-Trigger, gender & 0.479 & 0.454 & 0.452 & 0.680 & 0.734 & 0.444 & 0.660 \\
\addlinespace
Mistral & Gender-Secret & 0.406 & 0.440 & 0.414 & 0.400 & 0.765 & 0.497 & 0.784 \\
 & Soft-Trigger, greeting & 0.931 & 0.930 & 0.936 & 0.878 & 0.900 & 0.872 & 0.823 \\
 & Soft-Trigger, time & 0.929 & 0.946 & 0.949 & 0.275 & 0.930 & 0.875 & 0.858 \\
 & Soft-Trigger, gender & 0.948 & 0.935 & 0.940 & 0.901 & 0.892 & 0.873 & 0.828 \\
\addlinespace
Gemma & Soft-Trigger, greeting & 0.101 & 0.290 & 0.164 & 0.845 & 0.933 & 0.839 & 0.858 \\
 & Soft-Trigger, time & 0.916 & 0.719 & 0.876 & 1.000 & 0.962 & 0.860 & 0.875 \\
 & Soft-Trigger, gender & 0.560 & 0.804 & 0.763 & 0.794 & 0.951 & 0.849 & 0.868 \\
\bottomrule
\end{tabular}
}
\end{table}

\section*{Detection in the tail}

\Cref{tab:xm-operating} gives the share of lies each detector catches when the cutoff is set on half of a
dataset's own honest answers, at 1\% and 10\% of them flagged, averaged over the model's datasets, as
\cref{fig:k-operating} does for Llama.

\begin{table}[htbp]
\centering
\footnotesize
\caption{Share of lies caught at cutoffs set on each dataset's honest answers, at each model's mid block.}
\label{tab:xm-operating}
\fittable{% Generated by make_xmodel_assets.py from k_robustness.py operating points
\begin{tabular}{@{}l cc cc cc cc@{}}
\toprule
 & \multicolumn{2}{c}{Llama} & \multicolumn{2}{c}{Qwen$^\dagger$} & \multicolumn{2}{c}{Mistral} & \multicolumn{2}{c}{Gemma} \\
\cmidrule(lr){2-3}\cmidrule(lr){4-5}\cmidrule(lr){6-7}\cmidrule(lr){8-9}
Detector & 1\% & 10\% & 1\% & 10\% & 1\% & 10\% & 1\% & 10\% \\
\midrule
Base probe D2 & 0.32 & 0.66 & 0.02 & 0.17 & 0.13 & 0.47 & 0.28 & 0.66 \\
D5 & 0.34 & 0.58 & 0.02 & 0.12 & 0.13 & 0.47 & 0.08 & 0.30 \\
D9 & 0.47 & 0.77 & 0.02 & 0.14 & 0.14 & 0.49 & 0.43 & 0.78 \\
K & 0.21 & 0.56 & 0.02 & 0.13 & 0.21 & 0.54 & 0.07 & 0.26 \\
D9 + K & 0.38 & 0.73 & 0.03 & 0.13 & 0.23 & 0.59 & 0.30 & 0.63 \\
D5 + K & 0.32 & 0.67 & 0.02 & 0.11 & 0.22 & 0.59 & 0.24 & 0.53 \\
\bottomrule
\end{tabular}
}
\end{table}

\section*{The probe trained on the model's own lies}

\Cref{tab:xm-aisi} sets D3, the base probe trained on the model's own belief-labelled lies
(\cref{sec:supervised}), beside the probe trained on RepE and the three contrasts, for the two models whose
AISI rollouts we have. On both, at the middle block, D3 improves on the RepE probe and D9 stays ahead of
D3; at Gemma's block for layer~20, D3 leads.

\begin{table}[htbp]
\centering
\footnotesize
\caption{Mean AUROC over the model's datasets of D2, D3 and the contrasts, Llama and Gemma.}
\label{tab:xm-aisi}
\fittable{% Generated by make_xmodel_assets.py from aisi_probe.py compare
\begin{tabular}{@{}l ccc ccc@{}}
\toprule
 & \multicolumn{3}{c}{Llama} & \multicolumn{3}{c}{Gemma} \\
\cmidrule(lr){2-4}\cmidrule(lr){5-7}
Detector, at the block for layer & 20 & 40 & 44$^*$ & 20 & 40 & 44$^*$ \\
\midrule
Base probe D2, trained on RepE & 0.697 & 0.778 & 0.756 & 0.512 & 0.783 & 0.761 \\
D3, trained on the model's own lies & 0.782 & 0.878 & 0.872 & 0.802 & 0.858 & 0.770 \\
D5 & 0.750 & 0.798 & 0.791 & 0.575 & 0.679 & 0.691 \\
D8 & 0.863 & 0.894 & 0.880 & 0.622 & 0.484 & 0.564 \\
D9 & 0.887 & 0.899 & 0.884 & 0.591 & 0.900 & 0.801 \\
\bottomrule
\end{tabular}
}
\end{table}

\section*{How fully each model plays the roles}

\Cref{tab:xm-roleplay} gives, for each group of roles, the share of the model's rollouts the judge scored
fully in role. Gemma plays nearly every role fully, Llama and Mistral about two thirds of the time, and
Qwen under a third, least of all its deceptive and knowledge-denial roles.

\begin{table}[htbp]
\centering
\footnotesize
\caption{Share of each persona group's rollouts the judge scored 3, fully in role.}
\label{tab:xm-roleplay}
\fittable{% Generated by make_xmodel_assets.py from the judge's scores
\begin{tabular}{@{}l cccc@{}}
\toprule
Persona group & Llama & Qwen$^\dagger$ & Mistral & Gemma \\
\midrule
Deceptive (D5) & 58\% & 15\% & 73\% & 97\% \\
Honest (D5) & 74\% & 52\% & 85\% & 99\% \\
Knowledge-denial (D8) & 78\% & 10\% & 54\% & 96\% \\
Knowledge-candour (D8) & 54\% & 22\% & 53\% & 89\% \\
Sycophantic & 75\% & 37\% & 60\% & 97\% \\
Anti-sycophantic & 82\% & 28\% & 75\% & 100\% \\
\midrule All 35 roles & 69\% & 28\% & 68\% & 96\% \\
\bottomrule
\end{tabular}
}
\end{table}

\section*{The judge's filter}

\Cref{tab:xm-banks} sets the contrasts built from each model's strict bank, which keeps only rollouts the
judge scored 3, against those built from its lenient bank, which also keeps rollouts scored 2. Qwen has
only the lenient bank. On the other three models the two banks differ by at most \V{xm-bank-gap-max}
on any of these numbers.

\begin{table}[htbp]
\centering
\footnotesize
\caption{The contrasts from the strict and the lenient persona bank at each model's mid block.}
\label{tab:xm-banks}
\fittable{% Generated by make_xmodel_assets.py from bank_comparison
\begin{tabular}{@{}ll cccc@{}}
\toprule
 & & \multicolumn{3}{c}{Mean over the model's datasets} & \\
\cmidrule(lr){3-5}
Model & Bank & D5 & D8 & D9 & D9 transfer \\
\midrule
Llama & score 3 & 0.798 & 0.894 & 0.899 & 0.916 \\
 & scores 2--3 & 0.786 & 0.880 & 0.888 & 0.907 \\
Qwen$^\dagger$ & score 3 & n/a & n/a & n/a & n/a \\
 & scores 2--3 & 0.380 & 0.668 & 0.441 & 0.453 \\
Mistral & score 3 & 0.719 & 0.727 & 0.762 & 0.810 \\
 & scores 2--3 & 0.712 & 0.747 & 0.770 & 0.816 \\
Gemma & score 3 & 0.679 & 0.484 & 0.900 & 0.921 \\
 & scores 2--3 & 0.684 & 0.487 & 0.895 & 0.912 \\
\bottomrule
\end{tabular}
}
\end{table}

\section*{Other checks}

\Cref{tab:xm-checks} gives two checks. The first is the template-free K of \cref{sec:xm-k}: a plain system
turn (``You are a helpful assistant.''), the bare question, one of six instruction lines that ask whether
the model knows the answer and limit the reply to that, and the replies ``That is not something I know.''
and ``I have no knowledge of that.'', over each model's own known and unknowable questions. Its criteria
were fixed before any of its rows was read on the three further models; on Mistral it kept the
identical-text result but not the template K's transfer or mean. The second sets D9's mean token read
against the best of the six token reads of \cref{sec:token-read}.

\begin{table}[htbp]
\centering
\footnotesize
\caption{The template-free K beside the template K, and D9's token reads, at each model's mid block.}
\label{tab:xm-checks}
\fittable{% Generated by make_xmodel_assets.py from k_final and persona_positions
\begin{tabular}{@{}l ccc cc@{}}
\toprule
 & \multicolumn{3}{c}{Template-free K} & \multicolumn{2}{c}{D9's token read, mean AUROC} \\
\cmidrule(lr){2-4}\cmidrule(lr){5-6}
Model & Identical text & Transfer & Template K's transfer & Mean & Best \\
\midrule
Llama & 0.794 & 0.791 & 0.823 & 0.899 & 0.910 (maximum) \\
Qwen$^\dagger$ & n/a & 0.509 & 0.508 & 0.441 & 0.675 (first-token) \\
Mistral & 0.837 & 0.614 & 0.799 & 0.762 & 0.772 (top-10\% mean) \\
Gemma & n/a & 0.444 & 0.457 & 0.900 & 0.900 (mean) \\
\bottomrule
\end{tabular}
}
\end{table}

\FloatBarrier
